% Generated draft: requires mathematical and source review.
\documentclass[11pt]{article}
\usepackage[margin=25mm]{geometry}
\usepackage{fontspec}
\setmainfont{Latin Modern Roman}
\usepackage{amsmath,amssymb,mathtools,unicode-math}
\setmathfont{Latin Modern Math}
\usepackage{enumitem,needspace,xurl,hyperref}
\setlength{\parindent}{0pt}
\setlength{\parskip}{5pt}
\setlength{\emergencystretch}{3em}
\newcommand{\N}{\mathbb{N}}
\newcommand{\Z}{\mathbb{Z}}
\newcommand{\Q}{\mathbb{Q}}
\newcommand{\R}{\mathbb{R}}
\newcommand{\C}{\mathbb{C}}
\newcommand{\F}{\mathbb{F}}
\newcommand{\A}{\mathbb{A}}
\newcommand{\PP}{\mathbb{P}}
\newcommand{\E}{\mathbb{E}}
\newcommand{\eps}{\varepsilon}
\newcommand{\dd}{\,\mathrm{d}}
\begin{document}
\section*{UnsolvedMath ID: 30001547}
\textbf{Source problem code:} catalogue entry
\subsection{Definitions and mathematical statement}
The supplied question is: let \(L\) be the set of infinite words over an alphabet \(\Sigma\) avoiding a pattern \(p\); is \(L\) perfect? The record does not specify the alphabet, the pattern, the meaning of pattern avoidance, the topology, or the intended quantification over these parameters. Consequently, the following makes one conventional interpretation precise, conditionally; the supplied evidence does not establish that it is the intended formulation.

Under the one-sided interpretation, put \(\mathbb N_0=\{0,1,2,\ldots\}\) and regard an infinite word as a function \(w:\mathbb N_0\to\Sigma\). Give \(\Sigma\) the discrete topology and \(\Sigma^{\mathbb N_0}\) the product topology. For a finite word \(u\in\Sigma^*\), its prefix cylinder is
\[
[u]=\{w\in\Sigma^{\mathbb N_0}:w_0\cdots w_{|u|-1}=u\}.
\]
Here \(\Sigma^*\) denotes the finite words, including the empty word, and \(\Sigma^+\) denotes the nonempty finite words.

If a pattern means a nonempty finite word \(p\in V^+\) over a variable alphabet \(V\), and occurrences mean nonerasing morphic occurrences, define
\[
L_{\Sigma,p}=\left\{w\in\Sigma^{\mathbb N_0}:\begin{array}{l}
\text{there are no }i\in\mathbb N_0\text{ and monoid morphism }h:V^*\to\Sigma^*\\
\text{with }h(a)\ne\varepsilon\text{ for every }a\in V\text{ and}\\
w_iw_{i+1}\cdots w_{i+|h(p)|-1}=h(p)
\end{array}\right\}.
\]
The symbol \(\varepsilon\) denotes the empty word. This definition is conditional: the record does not say whether \(p\) is a variable pattern, a literal forbidden word, or another kind of pattern, nor whether erasing substitutions are permitted.

With these conventions, the question is whether \(L_{\Sigma,p}\) is perfect, meaning closed in \(\Sigma^{\mathbb N_0}\) with no isolated points. Equivalently, the requested property is the conjunction
\[
\overline{L_{\Sigma,p}}=L_{\Sigma,p}
\quad\text{and}\quad
\forall w\in L_{\Sigma,p}\ \forall n\in\mathbb N_0\ \exists v\in L_{\Sigma,p}:
\ v\ne w\ \text{ and }\ \forall j<n,\ v_j=w_j.
\]
Under this definition the empty set is perfect. Whether nonemptiness is intended as an additional requirement is unresolved. The source does not identify a particular \((\Sigma,p)\), state that the question concerns every such pair, or supply restrictions such as finiteness or cardinality of \(\Sigma\) or avoidability of \(p\). These omissions prevent identification of a unique fully quantified problem.
\subsection{Short English statement}
The question asks whether the space of infinite words avoiding \(p\) is perfect. Under a conventional product topology, this means that the space is closed and every allowed word has a different allowed word sharing any prescribed finite prefix. The alphabet, pattern interpretation, topology, and parameter restrictions remain unspecified, so the displayed formulation is conditional and requires mathematical review.
\subsection{Sources}
\begin{itemize}
\item Supplied local registry record: permanent UnsolvedMath ID \(30001547\), source problem code \texttt{catalogue entry}, title ``Perfection of Pattern-Avoiding Infinite-Word Spaces,'' in the Oberwolfach Reports Open Problems collection. No individual proposer is supplied; the proposed year is recorded as 2010.
\item Supplied report citation: \emph{Mini-Workshop: Combinatorics on Words (2011)}, Oberwolfach Reports, pp.~2195--2244, DOI \texttt{10.4171/owr/2010/37}. Supplied source URL and literature-source label ``original OWR report'': \url{https://doi.org/10.4171/owr/2010/37}. The citation's 2011 designation and the record's 2010 dating are retained without reconciliation.
\item Supplied literature-source label ``related repetition-free results'': \url{https://doczz.net/doc/8060441/repetition-free-words-yuji-kobayashi-1}. The record supplies no precise theorem or bibliographic attribution establishing its applicability to this question.
\item Literature assessment checked at \texttt{2026-08-21}: this is supplied dated curation, not an independently verified assessment. It labels the problem open, says that the report poses perfection of \(L\) as an open problem, and reports perfection results for some related repetition-free languages without identifying a general resolution of this extracted case. The background explicitly requests independent verification of time-sensitive status and cited claims.
\item The supplied statement status is \texttt{corrected\_verified}; its verification note says that the language, alphabet, and pattern symbols were restored to standard mathematical notation without changing the question. This records the supplied editorial qualification and does not establish mathematical completeness or independent source verification.
\end{itemize}
\end{document}
